\editionPage[\pdfbookmark[0]{Corrigés des défis}{bm-defi-cor-1-parent}]{defi-cor-1}{D01 · Corrigé}{Les nombres miroirs}{\hyperref[defi-ex-1]{Énoncé p.\,\pageref*{defi-ex-1}}}
\begin{blocEdition}{defi-cor-1}
\editionEtape{Fixer les bords}Le nombre commence par 4, donc il finit aussi par 4 ; le deuxième chiffre est égal au quatrième.\editionEtape{Compter les choix}Le deuxième chiffre a 10 valeurs possibles, le troisième aussi : $10\times10=\res{100}$ nombres miroirs.\editionEtape{Les extrêmes}Le plus petit est $\res{40\,004}$, le plus grand $\res{49\,994}$. Tous sont bien compris entre $40\,000$ et $50\,000$.
\end{blocEdition}
\editionPage{defi-cor-2}{D02 · Corrigé}{Somme et écart}{\hyperref[defi-ex-2]{Énoncé p.\,\pageref*{defi-ex-2}}}
\begin{blocEdition}{defi-cor-2}
\editionEtape{Égaliser les barres}On enlève l’écart au plus grand : il reste $1\,000-248=752$, soit deux barres égales.\editionEtape{Partager}Chaque barre vaut $752:2=376$ : le petit nombre est $\res{376}$ et le grand $376+248=\res{624}$.\editionEtape{Contrôler}$376+624=1\,000$ et $624-376=248$ : les deux conditions sont vérifiées.
\end{blocEdition}
\editionPage{defi-cor-3}{D03 · Corrigé}{Les élèves en rangs}{\hyperref[defi-ex-3]{Énoncé p.\,\pageref*{defi-ex-3}}}
\begin{blocEdition}{defi-cor-3}
\editionEtape{Un élève de plus}Avec un élève de plus, tous les rangs seraient complets : le nombre d’élèves plus un est un multiple commun de 4, 5 et 6.\editionEtape{Les multiples communs}Les multiples communs de 4, 5 et 6 sont 60, 120, 180, 240… Entre 101 et 201, on garde 120 et 180.\editionEtape{Conclure}Il y a $\res{119}$ ou $\res{179}$ élèves. Contrôle : $119=4\times29+3=5\times23+4=6\times19+5$.
\end{blocEdition}
\editionPage{defi-cor-4}{D04 · Corrigé}{Les divisions qui s’arrêtent}{\hyperref[defi-ex-4]{Énoncé p.\,\pageref*{defi-ex-4}}}
\begin{blocEdition}{defi-cor-4}
\editionEtape{Celles qui s’arrêtent}$1:2=\res{0{,}5}$ ; $1:4=\res{0{,}25}$ ; $1:5=\res{0{,}2}$ ; $1:8=\res{0{,}125}$.\editionEtape{Celles qui continuent}Pour 3 et 9, le reste 1 revient sans fin : $0{,}333\ldots$ et $0{,}111\ldots$ Pour 6, le reste 4 revient : $0{,}1666\ldots$ Pour 7, les six restes possibles tournent en boucle : $0{,}142857142857\ldots$\editionEtape{Retenir}Un reste qui réapparaît annonce une division sans fin : les mêmes chiffres se répètent.
\end{blocEdition}
\editionPage{defi-cor-5}{D05 · Corrigé}{La suite des lapins}{\hyperref[defi-ex-5]{Énoncé p.\,\pageref*{defi-ex-5}}}
\begin{blocEdition}{defi-cor-5}
\editionEtape{Poursuivre}$1 ;\ 1 ;\ 2 ;\ 3 ;\ 5 ;\ 8 ;\ 13 ;\ 21 ;\ 34 ;\ 55 ;\ 89 ;\ 144$ : le douzième terme est $\res{144}$.\editionEtape{Aller plus loin}Puis $233 ;\ 377 ;\ 610 ;\ 987 ;\ 1\,597$. Le terme $1\,597$ dépasse $1\,000$.\editionEtape{Compter}Les termes inférieurs à $1\,000$ sont les 16 premiers : $\res{16}$ termes.
\end{blocEdition}
\editionPage{defi-cor-6}{D06 · Corrigé}{La fourmi pressée}{\hyperref[defi-ex-6]{Énoncé p.\,\pageref*{defi-ex-6}}}
\begin{blocEdition}{defi-cor-6}
\editionEtape{Suivre les sauts}Positions : $\frac12$ ; $\frac34$ ; $\frac78$ ; $\frac{15}{16}$ ; $\res{\frac{31}{32}}$ de mètre.\editionEtape{Ce qui reste}Restes : $\frac12$, $\frac14$, $\frac18$, $\frac1{16}$, puis $\res{\frac1{32}}$ de mètre, à peine plus de 3 cm.\editionEtape{Observer}La fourmi s’approche du bout sans jamais l’atteindre par ces sauts : le reste est toujours divisé par deux, jamais nul.
\end{blocEdition}
\editionPage{defi-cor-7}{D07 · Corrigé}{Les douzièmes du milieu}{\hyperref[defi-ex-7]{Énoncé p.\,\pageref*{defi-ex-7}}}
\begin{blocEdition}{defi-cor-7}
\editionEtape{Convertir les bornes}$\frac13=\frac4{12}$ et $\frac34=\frac9{12}$.\editionEtape{Lister}Les numérateurs 5, 6, 7 et 8 conviennent : $\res{\frac5{12},\ \frac6{12},\ \frac7{12},\ \frac8{12}}$, soit 4 fractions.\editionEtape{Le demi}$\res{\frac6{12}=\frac12}$ : la moitié de l’unité est bien entre un tiers et trois quarts.
\end{blocEdition}
\editionPage{defi-cor-8}{D08 · Corrigé}{L’héritage des chameaux}{\hyperref[defi-ex-8]{Énoncé p.\,\pageref*{defi-ex-8}}}
\begin{blocEdition}{defi-cor-8}
\editionEtape{Avec 18 chameaux}La moitié : 9 ; le tiers : 6 ; le neuvième : 2. Total : $9+6+2=\res{17}$ chameaux : il en reste un, celui du voyageur.\editionEtape{Pourquoi}$\frac12+\frac13+\frac19=\frac9{18}+\frac6{18}+\frac2{18}=\frac{17}{18}$ : le testament ne partage pas tout le troupeau, il en laisse $\frac1{18}$.\editionEtape{Retenir}Chaque enfant reçoit un peu plus que prévu sur 17 chameaux, mais des nombres entiers : $\res{9,\ 6\text{ et }2}$.
\end{blocEdition}
\editionPage{defi-cor-9}{D09 · Corrigé}{Le décimal à code}{\hyperref[defi-ex-9]{Énoncé p.\,\pageref*{defi-ex-9}}}
\begin{blocEdition}{defi-cor-9}
\editionEtape{Traduire}Chiffres décimaux : $d$, puis 0, puis $2\times d$. Leur somme vaut $3\times d=9$.\editionEtape{Conclure}$d=3$ : les millièmes valent 6. Le nombre est $\res{3{,}306}$.\editionEtape{Contrôler}$6$ est le double de 3, le chiffre des centièmes est nul et $3+0+6=9$. Une seule valeur de $d$ convient.
\end{blocEdition}
\editionPage{defi-cor-10}{D10 · Corrigé}{La tirelire à deux pièces}{\hyperref[defi-ex-10]{Énoncé p.\,\pageref*{defi-ex-10}}}
\begin{blocEdition}{defi-cor-10}
\editionEtape{Un départ simple}Dix pièces de 20 c font 2 euros : il manque $1{,}50$ euro.\editionEtape{Échanger}Chaque échange ajoute $0{,}30$ euro : $1{,}50:0{,}30=5$ échanges.\editionEtape{Conclure}$\res5$ pièces de 50 c et $\res5$ pièces de 20 c. Contrôle : $5\times0{,}50+5\times0{,}20=2{,}50+1=3{,}50$ euros.
\end{blocEdition}
\editionPage{defi-cor-11}{D11 · Corrigé}{La balance aux trois masses}{\hyperref[defi-ex-11]{Énoncé p.\,\pageref*{defi-ex-11}}}
\begin{blocEdition}{defi-cor-11}
\editionEtape{Retrancher}Une masse posée avec l’objet compte en moins. Pour 2 g : objet $+1$ g d’un côté, 3 g de l’autre, car $\res{2=3-1}$.\editionEtape{D’autres exemples}5 g : $9-3-1=\res5$ ; 11 g : $9+3-1=\res{11}$ ; 13 g : $9+3+1$.\editionEtape{Tous les cas}En combinant chaque masse en plus, en moins ou absente, on obtient chaque entier de 1 à 13 : le tableau des 13 cas se vérifie un par un.
\end{blocEdition}
\editionPage{defi-cor-12}{D12 · Corrigé}{Le plus grand enclos}{\hyperref[defi-ex-12]{Énoncé p.\,\pageref*{defi-ex-12}}}
\begin{blocEdition}{defi-cor-12}
\editionEtape{Les rectangles}Longueur plus largeur valent 12 m : $11\times1$, $10\times2$, $9\times3$, $8\times4$, $7\times5$, $6\times6$.\editionEtape{Les aires}11, 20, 27, 32, 35 et 36 m$^2$.\editionEtape{Conclure}Le carré de 6 m de côté offre la plus grande aire : $\res{36\ \mathrm m^2}$. Plus les côtés sont proches, plus l’aire est grande.
\end{blocEdition}
\editionPage{defi-cor-13}{D13 · Corrigé}{L’horloge des angles}{\hyperref[defi-ex-13]{Énoncé p.\,\pageref*{defi-ex-13}}}
\begin{blocEdition}{defi-cor-13}
\editionEtape{Une heure du cadran}Trois heures font $90^\circ$, donc une heure fait $\res{30^\circ}$.\editionEtape{Les cas}À 2 h : $2\times30=\res{60^\circ}$ ; à 4 h : $4\times30=\res{120^\circ}$.\editionEtape{150 degrés}$150:30=5$ : à $\res{5\ \text h}$, et aussi à $\res{7\ \text h}$, où l’écart dans l’autre sens vaut 5 heures du cadran.
\end{blocEdition}
\editionPage{defi-cor-14}{D14 · Corrigé}{Les deux sabliers}{\hyperref[defi-ex-14]{Énoncé p.\,\pageref*{defi-ex-14}}}
\begin{blocEdition}{defi-cor-14}
\editionEtape{Départ}On retourne les deux sabliers. À 4 min, le petit est vide : on le retourne.\editionEtape{À 7 min}Le grand est vide : on le retourne. À 8 min, le petit est vide : le grand n’a coulé que depuis 1 min.\editionEtape{La dernière minute}On retourne alors le grand : il ne contient qu’une minute de sable en haut, qui s’écoule de 8 à $\res{9}$ minutes.
\end{blocEdition}
\editionPage{defi-cor-15}{D15 · Corrigé}{Le trésor à deux distances}{\hyperref[defi-ex-15]{Énoncé p.\,\pageref*{defi-ex-15}}}
\begin{blocEdition}{defi-cor-15}
\editionEtape{Deux cercles}Les points à 4 cm de P forment le cercle de centre P et de rayon 4 cm ; ceux à 3 cm de C, le cercle de centre C et de rayon 3 cm.\editionEtape{Les croisements}Comme $4+3=7>5$, les cercles se coupent en $\res2$ points, symétriques par rapport à la droite (PC).\editionEtape{Conclure}Le trésor est en l’un de ces deux points : il faudra un indice de plus pour choisir.
\end{blocEdition}
\editionPage{defi-cor-16}{D16 · Corrigé}{Les diagonales des polygones}{\hyperref[defi-ex-16]{Énoncé p.\,\pageref*{defi-ex-16}}}
\begin{blocEdition}{defi-cor-16}
\editionEtape{Depuis un sommet}Dans un hexagone, chaque sommet donne $6-3=3$ diagonales.\editionEtape{Éviter le double compte}$6\times3=18$ compte chaque diagonale deux fois : $\res9$ diagonales.\editionEtape{Les autres}Quadrilatère : $4\times1:2=\res2$ ; pentagone : $5\times2:2=\res5$.
\end{blocEdition}
\editionPage{defi-cor-17}{D17 · Corrigé}{Les mots au fil de l’eau}{\hyperref[defi-ex-17]{Énoncé p.\,\pageref*{defi-ex-17}}}
\begin{blocEdition}{defi-cor-17}
\editionEtape{Ce que fait le reflet}L’axe horizontal retourne chaque lettre de haut en bas ; l’ordre des lettres ne change pas, contrairement au miroir vertical.\editionEtape{CODE}C, O, D et E ont toutes un axe horizontal : le mot se lit $\res{\text{inchangé}}$.\editionEtape{Chercher}D’autres exemples : $\res{\text{BEC}}$, DÉCIDE, CHOC, ÉCHO. Un mot comme CAFÉ échoue : A et F n’ont pas d’axe horizontal.
\end{blocEdition}
\editionPage{defi-cor-18}{D18 · Corrigé}{Le grand cube peint}{\hyperref[defi-ex-18]{Énoncé p.\,\pageref*{defi-ex-18}}}
\begin{blocEdition}{defi-cor-18}
\editionEtape{Coins et arêtes}Les 8 coins ont $\res3$ faces peintes. Chaque arête porte 2 cubes non situés aux coins : $12\times2=\res{24}$ cubes à 2 faces.\editionEtape{Faces et intérieur}Chaque face a un carré central de $2\times2$ : $6\times4=\res{24}$ cubes à une face. L’intérieur forme un cube de $2\times2\times2=\res8$ cubes sans peinture.\editionEtape{Contrôler}$8+24+24+8=64$ : tous les cubes sont comptés.
\end{blocEdition}
\editionPage{defi-cor-19}{D19 · Corrigé}{Deux robinets ensemble}{\hyperref[defi-ex-19]{Énoncé p.\,\pageref*{defi-ex-19}}}
\begin{blocEdition}{defi-cor-19}
\editionEtape{Par minute}A : 5 fois moins de temps, 5 fois moins d’eau : $30:5=6$ L. B : $24:6=4$ L.\editionEtape{Ensemble}Chaque minute, $6+4=10$ L tombent dans la citerne.\editionEtape{Conclure}$270:10=\res{27}$ minutes. Contrôle : en 27 min, A donne 162 L et B 108 L, soit 270 L.
\end{blocEdition}
\editionPage{defi-cor-20}{D20 · Corrigé}{La récréation en secteurs}{\hyperref[defi-ex-20]{Énoncé p.\,\pageref*{defi-ex-20}}}
\begin{blocEdition}{defi-cor-20}
\editionEtape{Les effectifs}Ballon : $120:2=60$ ; corde : $120:3=40$ ; billes : $120-60-40=\res{20}$ élèves.\editionEtape{Le secteur}$\frac{20}{120}=\frac16$ : le secteur des billes occupe $\res{\frac16}$ du disque.\editionEtape{Contrôler}$\frac12+\frac13+\frac16=\frac36+\frac26+\frac16=1$ : le disque est entièrement partagé.
\end{blocEdition}
\editionPage{defi-cor-21}{D21 · Corrigé}{La somme la plus probable}{\hyperref[defi-ex-21]{Énoncé p.\,\pageref*{defi-ex-21}}}
\begin{blocEdition}{defi-cor-21}
\editionEtape{Le tableau}Les 36 cases sont équiprobables. La somme 7 apparaît sur une diagonale entière : 6 cases.\editionEtape{Le maximum}Aucune autre somme n’en a autant : $\res7$ est la plus probable, avec $\res{6\ \text{chances sur }36}$.\editionEtape{Les extrêmes}La somme 2 n’a qu’une case, (1 ; 1) ; la somme 12 aussi, (6 ; 6) : $\res{\text{mêmes chances}}$, 1 sur 36.
\end{blocEdition}
